\chapter{Marco teórico}

\section{Locomoción cuadrúpeda}

Un robot cuadrúpedo alterna fases de apoyo y oscilación para desplazar el cuerpo sin perder estabilidad. La secuencia temporal de estas fases define la marcha. En una marcha estática, el movimiento se ejecuta con suficiente lentitud para que el equilibrio pueda analizarse principalmente mediante la proyección del centro de masa y el polígono formado por los pies en contacto. En una marcha dinámica también deben considerarse la velocidad, la aceleración, las fuerzas de reacción y los efectos inerciales~\cite{zamani2011dynamics}.

Las plataformas cuadrúpedas modernas utilizan doce articulaciones actuadas, normalmente tres por extremidad. Esta configuración permite posicionar cada pie en tres dimensiones y controlar simultáneamente la postura del tronco~\cite{bledt2018cheetah3,hutter2016anymal}.

\section{Modelado cinemático}

La cinemática directa relaciona las posiciones articulares con la pose del pie respecto al cuerpo. Para una extremidad con tres grados de libertad puede expresarse de forma general como

\begin{equation}
    \mathbf{p}_{f}=f(\mathbf{q}),
    \label{eq:cinematica-directa}
\end{equation}

donde $\mathbf{q}$ contiene las posiciones articulares y $\mathbf{p}_{f}$ representa la posición cartesiana del pie. La cinemática inversa determina las posiciones articulares compatibles con una referencia cartesiana. Debido a los límites mecánicos y a la existencia de múltiples soluciones, el algoritmo debe comprobar el espacio alcanzable y seleccionar una configuración coherente con la geometría de la extremidad.

La relación diferencial entre velocidad articular y velocidad cartesiana se obtiene mediante el jacobiano:

\begin{equation}
    \dot{\mathbf{p}}_{f}=\mathbf{J}(\mathbf{q})\dot{\mathbf{q}}.
    \label{eq:jacobiano}
\end{equation}

Estas relaciones permiten generar trayectorias de los pies y transformarlas en referencias articulares para el controlador de posición.

\section{Dinámica y contacto}

El modelo dinámico de un cuadrúpedo con base flotante puede escribirse como

\begin{equation}
    \mathbf{M}(\mathbf{q})\ddot{\mathbf{q}}+
    \mathbf{h}(\mathbf{q},\dot{\mathbf{q}})=
    \mathbf{S}^{\mathsf{T}}\boldsymbol{\tau}+
    \mathbf{J}_{c}^{\mathsf{T}}\boldsymbol{\lambda},
    \label{eq:dinamica}
\end{equation}

donde $\mathbf{M}$ es la matriz de inercia, $\mathbf{h}$ agrupa los términos de Coriolis, centrífugos y gravitacionales, $\boldsymbol{\tau}$ contiene los esfuerzos actuados y $\boldsymbol{\lambda}$ representa las fuerzas de contacto. Los controladores de cuerpo completo y predictivos incorporan esta dinámica junto con restricciones de fricción, contacto y saturación~\cite{bellicoso2018dynamic,dicarlo2018mpc,neunert2018wholebody}.

\section{Control discreto de marcha}

Una máquina de estados finitos describe una marcha mediante posturas y transiciones. Cada estado identifica las extremidades en apoyo y oscilación, además de las referencias articulares o cartesianas. Para el gateo se mueve una extremidad a la vez con el propósito de conservar un triángulo de apoyo. Este enfoque proporciona una referencia nominal interpretable y permite incorporar límites antes de operar el hardware.

\section{Aprendizaje por refuerzo}

En aprendizaje por refuerzo, el agente observa el estado del sistema, ejecuta una acción y recibe una recompensa. El objetivo es aprender una política $\pi(\mathbf{a}\mid\mathbf{s})$ que maximice el retorno esperado. En locomoción, las observaciones pueden incluir orientación, velocidades, estado articular, contactos y fase de marcha; las acciones pueden corresponder a correcciones sobre referencias nominales.

Los resultados publicados muestran que la elección de la recompensa puede inducir patrones de locomoción coordinados y favorecer comportamientos energéticamente eficientes~\cite{heess2017emergence,fu2021energy}. En este proyecto la política no comandará directamente las señales PWM: generará correcciones acotadas que serán verificadas por el supervisor de seguridad.

\section{Simulación y transferencia}

La simulación permite estudiar contactos, perturbaciones y variaciones de parámetros de manera repetible. MuJoCo proporciona un motor físico orientado a sistemas articulados y control basado en modelos~\cite{todorov2012mujoco}. Para reducir la diferencia entre simulación y hardware se pueden variar masas, fricción, retardos, ruido y ganancias durante el entrenamiento. La transferencia se realizará progresivamente: simulación, ejecución sin actuadores, robot suspendido y pruebas sobre el suelo con soporte de seguridad.
